% Part: lambda-calculus
% Chapter: lambda-definability
% Section: partial-recursive-functions

\documentclass[../../../include/open-logic-section]{subfiles}

\begin{document}

\olfileid[pt]{lam}{ldf}{par}
\olsection{As funções recursivas parciais são \usetoken{S}{lambda definable}}

As funções recursivas parciais são aquelas obtidas a partir das funções
básicas por composição, recursão primitiva e minimização
ilimitada. Elas diferem das funções recursivas gerais porque não se
exige que as funções usadas na busca ilimitada sejam regulares.
Não exigir regularidade significa que, às vezes, as funções definidas por minimização
podem não estar definidas. 

À primeira vista, poderia parecer que os mesmos métodos usados para mostrar que
as funções recursivas gerais (totais) são todas !!{lambda definable} podem
ser usados para provar que todas as funções recursivas parciais são
!!{lambda definable}. Por exemplo, a composição de $f$ com $g$ é
!!{lambda define}d por $\lambd[x][F (G x)]$ se $f$ e $g$ são
!!{lambda define}d pelos termos $F$ e~$G$, respectivamente. No entanto, quando as
funções são parciais, isso é problemático. Quando $g(x)$ não está definida,
isso significa que $G x$ não tem forma normal. Na maioria dos casos, isso significa que $F (G
x)$ tampouco tem forma normal, que é o que queremos. Mas considere
o caso em que $F$ é $\lambd[x][\lambd[y][y]]$: nesse caso, $F (G x)$ de fato
tem uma forma normal ($\lambd[y][y]$). 

Esse problema não é insuperável, e há maneiras de
!!{lambda define} todas as funções recursivas parciais de modo que
valores indefinidos sejam representados por termos sem forma normal.
Essas maneiras, contudo, são um pouco mais complicadas e menos intuitivas
que a abordagem que adotamos para as funções recursivas gerais. Registramos
aqui o teorema sem prova:

\begin{thm}
  Todas as funções recursivas parciais são !!{lambda definable}.
\end{thm}

\end{document}
